\documentclass[10pt,a4paper]{amsart}
\usepackage[margin=19mm]{geometry}
\usepackage{amsmath,amssymb,mathtools}
\usepackage[scr=rsfs,cal=euler]{mathalfa}
\usepackage{xfrac}
\usepackage[unicode,hidelinks,bookmarks=false]{hyperref}
\newcommand{\ie}{i.e.\ }
\newcommand{\cf}{cf.\ }
\newcommand{\resp}{respectively\ }
\newcommand{\Subsection}{\subsection}
\providecommand{\nobreakdash}{\nobreak\hbox{-}}
\setlength{\emergencystretch}{2em}
\multlinegap0pt
\usepackage{amssymb}
\usepackage{enumerate}
\usepackage{xcolor}
\newtheorem{lemma}{Lemma}
\usepackage{cleveref}
\crefname{lemma}{Lemma}{Lemmas}
\newcommand{\N}{\mathbb{N}}
\newcommand{\R}{\mathbb{R}}
\newcommand{\les}{\lesssim}
\newcommand{\less}{\lessapprox}
\title{Weighted wave envelope estimates for the parabola}
\author{Jongchon Kim, Hyerim Ko}
\date{Source: arXiv:2511.04503v2 (CC BY 4.0)}
\begin{document}
\maketitle

\noindent\emph{Source and licence.} Weighted wave envelope estimates for the parabola. Version arXiv:2511.04503v2 (CC BY 4.0), distributed by arXiv under the Creative Commons Attribution 4.0 International licence, http://creativecommons.org/licenses/by/4.0/, as recorded in the arXiv licence metadata for this exact version and shown on the abstract page https://arxiv.org/abs/2511.04503v2. Full licence notice: \texttt{licenses/arxiv-2511.04503-CC-BY-4.0.txt}.\par\smallskip

\noindent\textit{Editorial scope.} Counted unit: the article's lemma bound (source0.tex lines 878--884). The article's complete original proof of it is delivered with it (source0.tex lines 1435--1472, one \texttt{proof} environment of 38 lines, which the article places away from the statement). No mathematics is written, repaired or re-derived: the statement and the proof are the article's own lines, in the article's own notation, and no retained line is substituted or re-flowed.  Retained uncounted as the source-local context the proof needs: lines 1420--1424.  Nothing was omitted from any retained span, so the package carries no omission record.  No retained line contains a citation, so the wrapper carries no bibliography and no bibliography entry.  No retained line contains a figure environment, an image-inclusion command or drawing code, so no image is delivered and the package declares no asset dependency.  Editorial formatting only, in addition: an AMS \texttt{amsart} preamble that restates the article's own declarations and macros used by the retained lines, namely the environments \texttt{lemma} on separate counters, together with the standard packages those lines need.  The retained environments print in this document as Lemma 1 in order of appearance, whereas the article prints the same items with the numbering of its own full document.  The complete native source of the article as distributed in the arXiv e-print for this exact version is bundled unchanged as \texttt{sources/188593/source0.tex}.
\begin{lemma}\label{lem:BG} Let $\{ a_i \}_{i\in I}$ be a sequence of non-negative real number indexed by a finite set $I$. For each $i\in I$, let $I_i \subset I$ be a subset containing $i$ such that $|I_i| \leq C_1$ for all $i\in I$ for some 
constant $C_1 \in \N$. Then there exists $C=C(C_1,p)$ such that for $p\ge1$,
\[ 
\Big(\sum_{i\in I} a_i\Big)^p \leq C\Big( \max_{i\in I} a_i^p + (\# I )^p\max_{\substack{i\in I,\\ j\notin I_i}} a_i^{\frac p2} a_j^{\frac p2}\Big). \]
\end{lemma}

\begin{lemma}\label{ptwise bound}
For any $x\in \R^2$, we have the pointwise bound
\begin{equation}\label{eqn:iteration}
  |f(x)|^p \less \sum_{|\theta|=R^{-1/2}} |f_{\theta}(x)|^{p} + \sum_{1\leq j\leq m} \sum_{|\tau|=K^{-j+1}} \sum_{\substack{\tau_1,\tau_2\subset \tau \\ |\tau_1|=|\tau_2| = K^{-j}  \\ d(\tau_1,\tau_2)\geq   K^{-j}}} |f_{\tau_1}(x)f_{\tau_2}(x)|^{\frac p2},
\end{equation}
where $d(\tau_1,\tau_2)$ denotes the distance between $\tau_1$ and $\tau_2$. 
\end{lemma}

\begin{proof}[Proof of \Cref{ptwise bound}]
Recall that $K\sim \log R$ is a dyadic number and $m$ is the integer such that 
\[1\le K\le \dots \le K^m \leq R^{1/2} < K^{m+1}.\]
Let $\mathcal T_0 =\{ \tau_0\}$ denote collection of the unit cube covering the parabola $\mathcal P$.

At the first stage, we decompose $\tau_0$ into a collection $\mathcal T_1=\mathcal T_1(\tau_0)$ of $K^{-1}\times K^{-2}$ boxes $\tau_1$ covering the $K^{-2}$-nbd of the parabola.
For each $\tau_1\in \mathcal T_1$,  let $\mathcal T_2(\tau_1)$ be a  collection of $K^{-2}\times K^{-4}$ boxes $\tau_2 \subset \tau_1$ covering $K^{-4}$-nbd of $\mathcal P$, and define $\mathcal T_2 = \cup_{\tau_1\in \mathcal T_1} \mathcal T_2(\tau_1)$.

Proceeding inductively, for $2\le j \le m$, we define $\mathcal T_j(\tau_{j-1})$ as the collection of boxes $\tau_j \subset \tau_{j-1}$ of dimension $K^{-j}\times K^{-2j}$ covering $K^{-2j}$-neighborhood of $\mathcal P$ and $\mathcal T_j$ similarly.
Finally, denote $\mathcal T_m=\{\theta\}$ and for each $j=1,\dots,m$ set
\[f_{\tau_j} = \sum_{\theta \subset \tau_j} f_\theta.\]


For each $\tau_1\in \mathcal T_1$, let 
\[
\mathcal N_1(\tau_1)=\{\tau_1'\in \mathcal T_1:~ \tau_1' \cap 2\tau_1 \neq \emptyset\}.
\]
It is clear that $\mathcal N_1(\tau_1)$ consists of only $O(1)$ many elements, and if $\tau_1' \notin \mathcal N_1(\tau_1)$ then $d(\tau_1,\tau_1')\ge 1/K$.
Since $\mathcal T_1$ is covered by such neighborhood $\mathcal N_1(\tau_1)$,
applying \Cref{lem:BG} to $|f|^p\leq (\sum_{\tau_1\in \mathcal T_1} |f_{\tau_1}|)^p$, and using that $\# \mathcal T_1 \les K$, we have
 \[ |f(x)|^p 
 \leq C\max_{\tau_1 \in \mathcal T_1} |f_{\tau_1}(x)|^p + CK^p \max_{\substack{\tau_1,\tau_1' \in \mathcal T_1; \\ d(\tau_1,\tau_1')\ge K^{-1} }} |f_{\tau_1}(x)|^{\frac p2} |f_{\tau_1'}(x)|^{\frac p2}   \]
for some absolute constant $C$.

Applying \Cref{lem:BG} again to the first term $|f_{\tau_1}|^p=|\sum_{\tau_2 \in \mathcal T_2(\tau_1)}f_{\tau_2}|^p$, we get 
\begin{align*}
  |f(x)|^p &\leq C^2 \max_{\tau_1} \max_{\tau_2\in \mathcal T_2(\tau_1)} |f_{\tau_2}(x)|^p + C^2K^p \max_{\tau_1} \max_{\substack{\tau_2,\tau_2' \in \mathcal T_2(\tau_1);\\ d(\tau_2,\tau_2')\ge K^{-2} }} |f_{\tau_2}(x)|^{\frac p2} |f_{\tau_2'}(x)|^{\frac p2} \\ &\qquad \quad+  CK^p \max_{\substack{\tau_1,\tau_1' \in \mathcal T_1;\\ d(\tau_1,\tau_1')\ge K^{-1} }} |f_{\tau_1}(x)|^{\frac p2} |f_{\tau_1'}(x)|^{\frac p2}.
\end{align*}
Continuing in this manner,  we get 
\begin{equation*}
	\begin{split}
  |f(x)|^p &\leq C^m \max_{\tau_m\in \mathcal{T}_m} |f_{\tau_m}(x)|^p + C^mK^p \sum_{j=1}^m \sum_{\tau_{j-1} \in \mathcal T_{j-1}} \max_{\substack{\tau_j,\tau_j' \in \mathcal T_j(\tau_{j-1});\\ d(\tau_j,\tau_j')\ge K^{-j} }} |f_{\tau_j}(x)|^{\frac p2} |f_{\tau_j'}(x)|^{\frac p2} \\
   &\less \sum_{|\theta|=R^{-1/2}} |f_{\theta}(x)|^{p} +  \sum_{j=1}^m \sum_{\tau_{j-1} \in \mathcal T_{j-1}} \sum_{\substack{\tau_j,\tau_j' \in \mathcal T_j(\tau_{j-1});\\ d(\tau_j,\tau_j')\ge K^{-j}  }} |f_{\tau_j}(x)|^{\frac p2} |f_{\tau_j'}(x)|^{\frac p2},
\end{split}
\end{equation*}
where we have used $C^m K^{O(1)} \less 1$ together with the fact that each $\tau_m \in \mathcal{T}_m$ contains $O(K)$ elements $\theta$ of size $|\theta| = R^{-1/2}$.
This completes the proof of Lemma \ref{ptwise bound}.
\end{proof}

\end{document}
